% ===================================================================
%  ਸਾਰਣੀ ਨੰ. 1 — ਸਕੂਲ। — ਹਾਈ ਸਕੂਲ। — ਜਮਾਤ।
%
%  Ceci n'est PAS une transcription : c'est une traduction, faite en
%  2026, en pendjabi oriental d'aujourd'hui, ecrit en GOURMOUKHI.
%  D'ou trois differences avec texto/io et texto/fr, et trois
%  seulement :
%
%   * pas de \nl ni de \cc — la traduction ne suit aucune ligne
%     imprimee, et les lignes de ce fichier ne sont que du confort ;
%   * pas de \begin{VUpage} — elle n'a ni feuillet ni folio, et la
%     page de lecture ne lui en affiche pas ;
%   * le gras se pose sur le TERME seul, ponctuation dehors.
%
%  Tout le reste est identique, et doit l'etre : les cles « %%K » sont
%  celles de texto/io, macro pour macro pour l'apparat, et les renvois
%  \textsuperscript{(n)} visent les MEMES objets de la planche — ce
%  sont eux qui ouvrent les gros plans.
%
%  CE N'EST PAS LA TRANSLITTERATION DE texto/pnb, ET C'EST LE POINT.
%  Le pendjabi s'ecrit en shahmoukhi au Pakistan et en gourmoukhi en
%  Inde ; les deux usages ont chacun leur lexique. La ou la colonne
%  « pnb » prend le mot persan ou arabe, celle-ci prend le mot indien,
%  parce que c'est celui qu'on imprime a Amritsar : la serie est une
%  « ਲੜੀ » et non un « سلسلہ », l'annexe est « ਸੰਲਗਨ » et non
%  « نال لگیا », le tableau noir est un « ਕਾਲਾ ਫੱਟਾ ». Une machine ne
%  fait pas ce choix-la ; il se fait mot par mot, et c'est pourquoi
%  les seize fichiers sont ecrits deux fois.
%
%  LES CHIFFRES RESTENT LATINS, renvois et numeros d'alinea : ils
%  numerotent la gravure et doivent se lire d'une colonne a l'autre.
%  C'est le parti du hindi, du bengali et des neuf autres.
%
%  LE PENDJABI POSTPOSE SES RAPPORTS et place le verbe a la fin : on
%  garde l'ordre des RENVOIS de l'ido, et l'on refait la phrase
%  autour.
%
%  LE NOM DE L'EDITEUR SE PRONONCE « del-MASS » : on ecrit ਦੇਲਮਾਸ,
%  non ਦੇਲਮਾ.
% ===================================================================

%%K t01-apar-0 apar
\VUsaut{2.0mm}
\VUcentre{12.6pt}{120}{ਵਿਆਖਿਆ-ਪੁਸਤਿਕਾ}
\VUsaut{3.2mm}
\VUcentre{8.4pt}{160}{ਸੰਲਗਨ}
\VUsaut{2.6mm}
\VUtitre{15.6pt}{0}{ਦੇਲਮਾਸ ਸਹਾਇਕ ਸਾਰਣੀਆਂ ਨਾਲ}
\VUsaut{3.4mm}
\VUornamento{9pt}
\VUsaut{5.0mm}
\VUcentre{13.2pt}{90}{ਪਹਿਲੀ ਲੜੀ}
\VUsaut{3.0mm}
\VUfilet{16mm}
\VUsaut{5.6mm}

%%K t01-apar-1 apar
\VUtitre{15.0pt}{60}{ਸਾਰਣੀ ਨੰ. 1}
\VUsaut{1.4mm}
\VUtitre{11.4pt}{0}{ਸਕੂਲ। --- ਹਾਈ ਸਕੂਲ। --- ਜਮਾਤ।}
\VUsaut{2.2mm}
\VUfilet{20mm}
\VUsaut{6.4mm}

%%K t01-tit-1 sub
\VUpk{10.2pt}{40}{ਆਮ ਵੇਰਵਾ।}
\VUsaut{3.0mm}

%%K t01-01-1 p
1. --- ਇਸ ਸਾਰਣੀ ਵਿਚ ਇਕ \VUgras{ਹਾਈ ਸਕੂਲ} ਦੀ \VUgras{ਜਮਾਤ} ਵਿਖਾਈ ਗਈ ਹੈ। ਇਸ ਜਮਾਤ
ਵਿਚ ਅੱਠ \VUgras{ਮੇਜ਼ਾਂ} \textsuperscript{(18)} ਤੇ ਅੱਠ \VUgras{ਬੈਂਚ}
\textsuperscript{(32)} ਹਨ। ਹਰ ਬੈਂਚ ਉੱਤੇ ਤਿੰਨ ਵਿਦਿਆਰਥੀ ਬਹਿੰਦੇ ਹਨ।
\VUgras{ਅਧਿਆਪਕ} \textsuperscript{(7)} \VUgras{ਕੁਰਸੀ} \textsuperscript{(8)} ਉੱਤੇ
ਆਪਣੇ \VUgras{ਮੰਚ} \textsuperscript{(6)} ਦੇ ਪਿੱਛੇ ਬੈਠੇ ਹਨ।

%%K t01-02-1 p
2. --- ਮੰਚ ਦੇ ਖੱਬੇ ਪਾਸੇ ਸ਼ੀਸ਼ੇ ਦੀ ਇਕ \VUgras{ਅਲਮਾਰੀ} \textsuperscript{(15)} ਖੜੀ
ਹੈ। ਸੱਜੇ ਪਾਸੇ \VUgras{ਕਾਲਾ ਫੱਟਾ} \textsuperscript{(1)} ਹੈ, ਤੇ ਉਸ ਉੱਤੇ
\VUgras{ਝਾੜਨ} \textsuperscript{(2)} ਤੇ \VUgras{ਚਾਕ} \textsuperscript{(3)}।

%%K t01-02-2 p
ਮੰਚ ਦੇ ਉੱਤੇ \VUgras{ਕੰਧ-ਚਿੱਤਰ} \textsuperscript{(9, 11, 12)} ਤੇ ਇਕ
\VUgras{ਨਕਸ਼ਾ} \textsuperscript{(10)} ਟੰਗੇ ਹੋਏ ਹਨ।

%%K t01-03-1 p
3. --- ਸਾਰੇ ਵਿਦਿਆਰਥੀ ਆਪਣੀ \VUgras{ਥਾਂ} \textsuperscript{(17)} ਉੱਤੇ ਨਹੀਂ। ਜਾਨ
\textsuperscript{(4)} ਕਾਲੇ ਫੱਟੇ ਕੋਲ ਖੜਾ ਹੈ; ਉਹ ਝਾੜਨ ਫੜੀ ਹੋਈ ਹੈ ਤੇ
\VUgras{ਰੇਖਾ-ਗਣਿਤ ਦੀਆਂ ਸ਼ਕਲਾਂ} ਮਿਟਾ ਰਿਹਾ ਹੈ।

%%K t01-03-2 p
ਪੀਟਰ ਮੰਚ ਕੋਲ ਹੈ; ਉਹ \VUgras{ਲਿਖਤੀ ਅਭਿਆਸ} \textsuperscript{(5)} ਇਕੱਠੇ ਕਰ ਕੇ
ਅਧਿਆਪਕ ਕੋਲ ਲੈ ਜਾ ਰਿਹਾ ਹੈ।

%%K t01-04-1 p
4. --- ਇਹ ਅਧਿਆਪਕ \VUgras{ਆਧੁਨਿਕ ਭਾਸ਼ਾਵਾਂ} ਦੇ ਅਧਿਆਪਕ ਹਨ। ਉਹ ਵਿਦਿਆਰਥੀਆਂ ਨੂੰ
ਪਰਾਈਆਂ ਭਾਸ਼ਾਵਾਂ \VUgras{(ਜਰਮਨ, ਅੰਗਰੇਜ਼ੀ, ਸਪੇਨੀ, ਇਤਾਲਵੀ, ਰੂਸੀ} ਜਾਂ
\VUgras{ਫ਼ਰਾਂਸੀਸੀ)} ਬੋਲਣਾ, ਪੜ੍ਹਨਾ ਤੇ ਲਿਖਣਾ ਸਿਖਾਉਂਦੇ ਹਨ।

%%K t01-05-1 p
5. --- ਜਮਾਤ ਬਹੁਤ ਵੱਡੀ ਤੇ ਬਹੁਤ ਰੌਸ਼ਨ ਹੈ। ਸਿਰੇ ਉੱਤੇ ਹਵਾ ਤੇ ਚਾਨਣ ਆਉਣ ਲਈ ਦੋ
\VUgras{ਰੌਸ਼ਨਦਾਨਾਂ} \textsuperscript{(78)} ਵਾਲੀ ਇਕ ਚੌੜੀ \VUgras{ਖਿੜਕੀ} ਤੇ ਇਕ
\VUgras{ਸ਼ੀਸ਼ੇ ਦਾ ਦਰਵਾਜ਼ਾ} ਹੈ।

%%K t01-05-2 p
ਇਹ ਜਮਾਤ ਠੰਢੀ ਨਹੀਂ। ਵਿਚਕਾਰ ਉਸ ਨੂੰ ਨਿੱਘੀ ਰੱਖਣ ਲਈ ਇਕ ਬਹੁਤ ਚੰਗੀ
\VUgras{ਅੰਗੀਠੀ} \textsuperscript{(46)} ਰੱਖੀ ਹੈ, ਆਪਣੀ \VUgras{ਨਾਲੀ}
\textsuperscript{(45)} ਸਮੇਤ।

%%K t01-05-3 p
ਕੰਧ ਉੱਤੇ \VUgras{ਖੂੰਟੀਆਂ} \textsuperscript{(58)} ਗੱਡੀਆਂ ਹੋਈਆਂ ਹਨ; ਉਨ੍ਹਾਂ ਉੱਤੇ
ਵਿਦਿਆਰਥੀਆਂ ਦੀਆਂ ਟੋਪੀਆਂ ਤੇ ਕੋਟ ਟੰਗੇ ਰਹਿੰਦੇ ਹਨ।

%%K t01-tit-2 sub
\VUsaut{5.2mm}
\VUpk{10.2pt}{40}{ਖੇਡ ਦਾ ਵਿਹੜਾ।}
\VUsaut{3.6mm}

%%K t01-06-1 p
6. --- \VUgras{ਘੜੀ} \textsuperscript{(63)} ਨੌਂ ਵੱਜ ਕੇ ਪੰਜ ਮਿੰਟ ਵਿਖਾ ਰਹੀ ਹੈ।

%%K t01-06-2 p
\VUgras{ਛੁੱਟੀ} \textsuperscript{(76)} ਹੁਣੇ ਮੁੱਕੀ ਹੈ ਤੇ ਪਾਠ ਫਿਰ ਸ਼ੁਰੂ ਹੋ ਰਿਹਾ
ਹੈ। ਛੁੱਟੀ \VUgras{ਵਿਹੜੇ} \textsuperscript{(75)} ਵਿਚ ਹੁੰਦੀ ਹੈ। ਸਾਨੂੰ ਕੁਝ
ਵਿਦਿਆਰਥੀ ਖੇਡਦੇ ਵਿਖਾਈ ਦਿੰਦੇ ਹਨ। ਇਕ \VUgras{ਸਹਾਇਕ ਅਧਿਆਪਕ} \textsuperscript{(74)}
ਉਨ੍ਹਾਂ ਉੱਤੇ ਨਜ਼ਰ ਰੱਖੀ ਹੋਈ ਹੈ।

%%K t01-07-1 p
7. --- ਵਿਹੜੇ ਵਿਚ ਸਾਨੂੰ ਹੋਰ ਵੀ ਕਈ ਲੋਕ ਵਿਖਾਈ ਦਿੰਦੇ ਹਨ : \VUgras{ਨਿਰੀਖਕ}
\textsuperscript{(73)}, \VUgras{ਮੁੱਖ ਅਧਿਆਪਕ} \textsuperscript{(71)} ਤੇ
\VUgras{ਅਨੁਸ਼ਾਸਕ} \textsuperscript{(72)}।

%%K t01-07-2 p
ਨਿਰੀਖਕ ਸਕੂਲਾਂ ਦਾ ਨਿਰੀਖਣ ਕਰਦੇ ਹਨ, ਮੁੱਖ ਅਧਿਆਪਕ ਸਕੂਲ ਚਲਾਉਂਦੇ ਹਨ, ਅਨੁਸ਼ਾਸਕ
ਵਿਦਿਆਰਥੀਆਂ ਦੀ ਪੜ੍ਹਾਈ ਉੱਤੇ ਨਜ਼ਰ ਰੱਖਦੇ ਹਨ।

%%K t01-07-3 p
ਚੌਥੇ ਬੰਦੇ \VUgras{ਚਪੜਾਸੀ} \textsuperscript{(77)} ਹਨ। ਉਹ ਗ਼ੈਰ-ਹਾਜ਼ਰ ਵਿਦਿਆਰਥੀਆਂ
ਦੇ ਨਾਂ ਲਿਖਣ ਲਈ ਕੱਛ ਵਿਚ ਇਕ ਰਜਿਸਟਰ ਦੱਬੀ ਹੋਏ ਹਨ।

%%K t01-08-1 p
8. --- ਵਿਹੜੇ ਦੇ ਸਿਰੇ ਉੱਤੇ \VUgras{ਕਸਰਤ ਦੇ ਸਾਧਨ} \textsuperscript{(70)} ਤੇ
\VUgras{ਦਸਤਕਾਰੀ} \textsuperscript{(69)} ਦੀ ਕਾਰਜਸ਼ਾਲਾ ਹੈ।

%%K t01-08-2 p
ਸਾਰਣੀ ਦੇ ਸੱਜੇ ਪਾਸੇ \VUgras{ਸੰਗੀਤ} ਤੇ \VUgras{ਚਿੱਤਰਕਲਾ} ਦੇ \VUgras{ਕਮਰੇ}
ਮੁਸ਼ਕਲ ਨਾਲ ਵਿਖਾਈ ਦਿੰਦੇ ਹਨ।

%%K t01-noto-1 noto
\VUnotes{143.2mm}{%
(*) \textit{ਨਾਵਾਂ} ਲਈ ਵੀ ਇਹੋ ਸਲਾਹ ਦਿੱਤੀ ਜਾਂਦੀ ਹੈ ਕਿ ਉਨ੍ਹਾਂ ਨੂੰ ਉਨ੍ਹਾਂ ਦੇ
ਲਾਤੀਨੀ ਰੂਪ ਵਿਚ ਲਿਖਿਆ ਜਾਵੇ। ਅਣਗਿਣਤ ਰੂਪਾਂ ਤੋਂ ਬਚਣ ਦਾ ਇਹੋ ਇਕ ਢੰਗ ਹੈ। ਫਿਰ,
\textit{Giovanni, Jean, Johann, Jan, Hans, John, Ivan} ਵਿਚੋਂ ਚੁਣੀਏ ਵੀ ਕਿਵੇਂ?
ਕੀ ਨਾਵਾਂ ਦਾ ਵੱਖਰਾ ਕੋਸ਼ ਬਣਾਇਆ ਜਾਵੇ? ਲਾਤੀਨੀ ਕਰਤਾ-ਰੂਪ ਲੈ ਲਵੋ ਤੇ ਆਖੋ :
\VUgras{Ioannes, Ioanna; Iulius, Iulia; Iakobus; Andreas; Lukas.}
(ਸੰਪੂਰਨ ਵਿਆਕਰਣ)।%
}

%%K t01-tit-3 sub
\VUpk{10.2pt}{40}{ਵਿਸਤਾਰ ਨਾਲ ਵੇਰਵਾ।}
\VUsaut{2.4mm}

%%K t01-tit-4 sub
\VUpk{10.2pt}{40}{ਚੰਗੇ ਵਿਦਿਆਰਥੀ।}
\VUsaut{3.0mm}

%%K t01-09-1 p
9. --- ਮੰਚ ਦੇ ਸੱਜੇ ਪਾਸੇ ਪਹਿਲੇ ਬੈਂਚ ਉੱਤੇ ਬੈਠੇ ਵਿਦਿਆਰਥੀ \VUgras{ਹੈਨਰੀ} (*),
\VUgras{ਸਟੀਫ਼ਨ} ਤੇ \VUgras{ਚਾਰਲਸ} ਹਨ।

%%K t01-09-2 p
ਹੈਨਰੀ ਬਹੁਤ ਧਿਆਨ ਦੇਣ ਵਾਲਾ ਤੇ ਮਿਹਨਤੀ ਹੈ; ਉਹ ਅਧਿਆਪਕ ਦੀ ਗੱਲ ਸੁਣ ਰਿਹਾ ਹੈ। ਉਸ ਦੀ
ਮੇਜ਼ ਉੱਤੇ ਇਕ \VUgras{ਕਾਪੀ} \textsuperscript{(28)}, ਇਕ \VUgras{ਕਿਤਾਬ}
\textsuperscript{(29)}, ਇਕ \VUgras{ਸ਼ਬਦਕੋਸ਼} \textsuperscript{(30)} ਤੇ ਇਕ
\VUgras{ਕਲਮਦਾਨ} \textsuperscript{(31)} ਹੈ। ਉਸ ਦੀ ਕਿਤਾਬ ਵਿਚ ਬਹੁਤ ਸਾਰੇ
\VUgras{ਸਫ਼ੇ} ਹਨ; ਹਰ ਅਧਿਆਇ ਵਿਚ ਕਈ \VUgras{ਪੈਰੇ} ਹਨ ਤੇ ਹਰ \VUgras{ਸਤਰ} ਵਿਚ ਬਹੁਤ
ਸਾਰੇ \VUgras{ਸ਼ਬਦ}।

%%K t01-09-4 p
ਉਸ ਦਾ ਗੁਆਂਢੀ ਸਟੀਫ਼ਨ \textsuperscript{(24)} ਜੋਸ਼ੀਲਾ ਹੈ। ਉਹ ਅਗਲੀ ਵਾਰ ਦਾ ਦਿੱਤਾ ਪਾਠ
ਆਪਣੀ ਕਾਪੀ ਵਿਚ ਲਿਖ ਰਿਹਾ ਹੈ। ਉਸ ਦੀ \VUgras{ਲਿਖਾਈ} ਸੋਹਣੀ ਹੈ। ਉਸ ਦੀ ਕਿਤਾਬ ਬੰਦ
ਹੈ। ਸਿਰਫ਼ \VUgras{ਜਿਲਦ} \textsuperscript{(27)} ਵਿਖਾਈ ਦਿੰਦੀ ਹੈ। ਉਸ ਕੋਲ ਉਸ ਦਾ
\VUgras{ਪਰਕਾਰ} \textsuperscript{(26)} ਪਿਆ ਹੈ।

%%K t01-09-5 p
ਚਾਰਲਸ \textsuperscript{(23)} ਹੱਥ ਵਿਚ ਕਿਤਾਬ ਫੜੀ ਹੋਈ ਹੈ; ਉਹ ਉਸ ਨੂੰ ਖੋਲ੍ਹ ਕੇ ਆਪਣਾ
ਪਾਠ ਫਿਰ ਦੁਹਰਾ ਰਿਹਾ ਹੈ। ਹਰ ਵਿਦਿਆਰਥੀ ਵਾਂਗ ਉਸ ਦੇ ਸਾਹਮਣੇ ਵੀ ਇਕ \VUgras{ਦਵਾਤ}
\textsuperscript{(25)} ਹੈ। ਉਹ ਆਗਿਆਕਾਰ ਹੈ।

%%K t01-10-1 p
10. --- ਅਗਲੀ ਮੇਜ਼ ਉੱਤੇ \VUgras{ਲੂਈ, ਐਂਟਨੀ} ਤੇ \VUgras{ਪਾਲ} ਹਨ। ਲੂਈ ਆਪਣਾ
\VUgras{ਪਾਠ} \textsuperscript{(22)} ਸੁਣਾ ਰਿਹਾ ਹੈ ਤੇ ਅਧਿਆਪਕ ਦੇ \VUgras{ਸਵਾਲਾਂ}
ਦਾ ਉੱਤਰ ਦੇ ਰਿਹਾ ਹੈ। ਐਂਟਨੀ \textsuperscript{(21)} ਬਹੁਤ ਬੇਧਿਆਨ ਹੈ; ਅਧਿਆਪਕ ਕਦੇ
ਕਦੇ ਉਸ ਨੂੰ ਸਜ਼ਾ ਵੀ ਦਿੰਦੇ ਹਨ। ਪਾਲ \textsuperscript{(20)} ਚੰਗਾ ਮਿੱਤਰ ਹੈ, ਮਿਲਣਸਾਰ
ਤੇ ਉਪਕਾਰੀ। ਉਹ ਬਹੁਤ ਧਿਆਨ ਦਿੰਦਾ ਹੈ।

%%K t01-tit-5 sub
\VUsaut{4.2mm}
\VUpk{10.2pt}{40}{ਮਾੜੇ ਵਿਦਿਆਰਥੀ।}
\VUsaut{3.2mm}

%%K t01-11-1 p
11. --- ਹੈਨਰੀ ਦੇ ਪਿੱਛੇ \VUgras{ਜਾਰਜ} \textsuperscript{(38)} ਬੈਠਾ ਹੈ। ਉਹ ਮਾੜਾ
ਮੁੰਡਾ ਨਹੀਂ, ਪਰ \VUgras{ਗਾਲੜੀ}, ਬੇਧਿਆਨ ਤੇ ਸ਼ਰਾਰਤੀ ਹੈ। ਉਸ ਦੀ \VUgras{ਚੰਚਲਤਾ} ਤੇ
\VUgras{ਅਵੱਗਿਆ} ਪਾਠ ਦੇ ਵੇਲੇ \VUgras{ਵਿਘਨ} ਪਾਉਂਦੀਆਂ ਹਨ, ਤੇ ਉਸ ਨੂੰ ਅਕਸਰ ਸਖ਼ਤ
\VUgras{ਚੇਤਾਵਨੀ} ਮਿਲਣੀ ਚਾਹੀਦੀ ਹੈ। ਉਸ ਦੀ \VUgras{ਕੱਚੀ ਕਾਪੀ} \textsuperscript{(35)}
ਗੰਦੀ ਹੈ, ਉਹ ਉਸ ਉੱਤੇ ਆਪਣੀ \VUgras{ਕਲਮ} \textsuperscript{(34)} ਨਾਲ
\VUgras{ਸਿਆਹੀ ਦੇ ਧੱਬੇ} ਪਾਉਂਦਾ ਰਹਿੰਦਾ ਹੈ, ਇਸੇ ਲਈ ਉਸ ਨੂੰ ਹਮੇਸ਼ਾ ਆਪਣਾ
\VUgras{ਰਬੜ} \textsuperscript{(36)} ਚਾਹੀਦਾ ਹੈ; ਉਸ ਦਾ \VUgras{ਬਸਤਾ}
\textsuperscript{(33)} ਪਾਟਿਆ ਹੋਇਆ ਹੈ, ਕਿਉਂਕਿ ਉਹ ਬਹੁਤ ਲਾਪਰਵਾਹ ਹੈ। ਆਧੁਨਿਕ
ਭਾਸ਼ਾਵਾਂ ਦੀ ਜਮਾਤ ਵਿਚ ਉਹ ਆਪਣਾ \VUgras{ਪੈਨਸਿਲ-ਧਾਰਕ} \textsuperscript{(37)} ਕਿਉਂ
ਲਿਆਉਂਦਾ ਹੈ?

%%K t01-11-2 p
ਉਸ ਦਾ ਗੁਆਂਢੀ ਐਡਮੰਡ \textsuperscript{(39)} ਉਸ ਨੂੰ ਪਸੰਦ ਨਹੀਂ ਕਰਦਾ। ਉਹ ਮੰਨੀਏ ਤਾਂ
ਥੋੜਾ ਆਲਸੀ ਹੈ, ਪਰ ਚੁੱਪ-ਚੁਪੀਤਾ ਤੇ ਟਿਕਿਆ ਹੋਇਆ ਹੈ। ਉਹ ਬਕਬਕ ਨਹੀਂ ਕਰਦਾ; ਬੱਸ,
ਸੁਣਨ ਦੀ ਥਾਂ ਉਹ ਆਪਣੀ \VUgras{ਪਾਠ-ਪੁਸਤਕ} ਦੀਆਂ \VUgras{ਕਹਾਣੀਆਂ} ਪੜ੍ਹਨਾ ਵੱਧ ਪਸੰਦ
ਕਰਦਾ ਹੈ।

%%K t01-11-3 p
ਇਸ ਬੈਂਚ ਦਾ ਤੀਜਾ ਵਿਦਿਆਰਥੀ \VUgras{ਲਡੋਵਿਕ} \textsuperscript{(40)} ਹੈ। ਉਹ ਮਿਹਨਤੀ
ਹੈ। ਉਹ \VUgras{ਕਾਗਜ਼ ਦੇ} ਇਕ ਚਿੱਟੇ \VUgras{ਸਫ਼ੇ} ਉੱਤੇ ਲਿਖ ਰਿਹਾ ਹੈ। ਆਪਣੇ ਸਾਹਮਣੇ
ਉਸ ਨੇ ਆਪਣਾ \VUgras{ਸੋਖ਼ਤਾ} \textsuperscript{(41)} ਰੱਖਿਆ ਹੋਇਆ ਹੈ।

%%K t01-12-1 p
12. --- ਅਧਿਆਪਕ ਦੇ ਖੱਬੇ ਪਾਸੇ ਦੂਜੇ ਬੈਂਚ ਦੇ ਵਿਦਿਆਰਥੀ ਬਹੁਤ ਛੋਟੇ ਹਨ। ਪਹਿਲਾ, ਨਿੱਕਾ
\VUgras{ਐਮਿਲ}, ਹਾਲੇ ਚੰਗੀ ਤਰ੍ਹਾਂ ਲਿਖ ਨਹੀਂ ਸਕਦਾ; ਉਹ ਆਪਣੀ \VUgras{ਸਲੇਟ}
\textsuperscript{(42)}, ਆਪਣੀ \VUgras{ਸਲੇਟੀ} ਤੇ ਆਪਣਾ \VUgras{ਸਪੰਜ}
\textsuperscript{(43)} ਲਿਆਉਂਦਾ ਹੈ।

%%K t01-12-2 p
ਉਸ ਕੋਲ \VUgras{ਔਗਸਤ} ਤੇ \VUgras{ਬਰਨਾਰਡ} \textsuperscript{(44)} ਹਨ। ਪਿਛਲਾ ਚੰਗਾ
ਕੰਮ ਕਰਦਾ ਹੈ, ਪਰ ਉਸ ਦਾ \VUgras{ਪਹਿਰਾਵਾ} ਬਹੁਤ ਮੈਲਾ-ਕੁਚੈਲਾ ਰਹਿੰਦਾ ਹੈ।

%%K t01-13-1 p
13. --- ਉਨ੍ਹਾਂ ਦੇ ਪਿੱਛੇ ਤਿੰਨ \VUgras{ਆਲਸੀ} ਬੈਠੇ ਹਨ। \VUgras{ਐਲਬਰਟ} ਆਪਣੇ ਪਾਠ
ਯਾਦ ਨਹੀਂ ਕਰਦਾ। \VUgras{ਜੇਮਜ਼} \textsuperscript{(47)} ਸੁਣਨ ਦੀ ਥਾਂ ਸੁੱਤਾ ਪਿਆ ਹੈ।
ਉਸ ਦਾ \VUgras{ਆਲਸ} ਉਸ ਨੂੰ \VUgras{ਝਿੜਕ} ਤੇ \VUgras{ਸਜ਼ਾ} ਤਕ ਵੀ ਪਹੁੰਚਾਉਂਦਾ ਹੈ।
ਅੰਤ ਵਿਚ \VUgras{ਕ੍ਰਿਸਚੀਅਨ} \textsuperscript{(48)} ਕਿਤੇ ਵੱਧ ਸ਼ਰਾਰਤੀ ਹੈ। ਉਸ ਦਾ
\VUgras{ਆਚਰਣ} ਮਾੜਾ ਹੈ ਤੇ ਉਹ ਸੁਧਰਨ ਦਾ ਕੋਈ ਯਤਨ ਨਹੀਂ ਕਰਦਾ।

%%K t01-14-1 p
14. --- ਉਨ੍ਹਾਂ ਦੇ ਗੁਆਂਢੀ, ਇਸ ਦੇ ਉਲਟ, ਸੁਸ਼ੀਲ ਹਨ। \VUgras{ਅਰਨੈਸਟ}
\textsuperscript{(51)} ਨੂੰ ਤਰਤੀਬ ਤੇ ਨੇਮ ਪਿਆਰੇ ਹਨ; ਉਹ ਸਦਾ ਆਪਣਾ \VUgras{ਫੁੱਟਾ}
\textsuperscript{(50)} ਤੇ ਆਪਣੀ \VUgras{ਪੈਨਸਿਲ} \textsuperscript{(49)} ਕੰਮ ਵਿਚ
ਲਿਆਉਂਦਾ ਹੈ। ਉਹ \VUgras{ਬਾਹਰਲਾ ਵਿਦਿਆਰਥੀ} ਹੈ।

%%K t01-14-2 p
\VUgras{ਫ਼ਰਾਂਸਿਸ} \textsuperscript{(52)} \VUgras{ਹੋਸਟਲੀਆ} ਹੈ; ਉਹ ਵਰਦੀ ਪਾਉਂਦਾ
ਹੈ, ਤੇ ਸਕੂਲ ਵਿਚ ਹੀ ਖਾਂਦਾ ਤੇ ਸੌਂਦਾ ਹੈ। ਉਹ ਚੰਗਾ \VUgras{ਮਿੱਤਰ} ਹੈ।

%%K t01-14-3 p
ਮੌਰਿਸ ਆਪਣੇ \VUgras{ਚਾਕੂ} \textsuperscript{(56)} ਨਾਲ ਆਪਣੀ \VUgras{ਪੈਨਸਿਲ} ਛਿੱਲ
ਰਿਹਾ ਹੈ; ਉਹ ਬਹੁਤ ਸਾਵਧਾਨ ਹੈ।

%%K t01-14-4 p
ਉਹ ਆਪਣੀਆਂ ਸਾਰੀਆਂ ਕਿਤਾਬਾਂ ਲਿਆਉਂਦਾ ਹੈ, ਆਪਣਾ \VUgras{ਵਿਆਕਰਣ}
\textsuperscript{(55)}, ਆਪਣੀ \VUgras{ਨੋਟਬੁੱਕ} \textsuperscript{(54)}, ਤੇ ਇਕ
\VUgras{ਲਿਖਣ ਦਾ ਗੱਦਾ} \textsuperscript{(53)} ਵੀ। ਉਸ ਦਾ \VUgras{ਥੈਲਾ}
\textsuperscript{(57)} ਉਸ ਦੇ ਪਿੱਛੇ ਭੁੰਜੇ ਪਿਆ ਹੈ।

%%K t01-14-5 p
ਜਮਾਤ ਦੇ ਪਿਛਲੇ ਪਾਸੇ ਦੀਆਂ ਦੋ ਮੇਜ਼ਾਂ ਉੱਤੇ \VUgras{ਚੰਗੇ ਵਿਦਿਆਰਥੀ}
\textsuperscript{(19)} ਬੈਠੇ ਹਨ।

%%K t01-tit-6 sub
\VUsaut{4.6mm}
\VUpk{10.2pt}{40}{ਚਿੱਤਰਕਲਾ ਤੇ ਸੰਗੀਤ।}
\VUsaut{3.4mm}

%%K t01-15-1 p
15. --- \VUgras{ਚਿੱਤਰਕਲਾ-ਕਮਰੇ} ਵਿਚ ਕੰਧ ਉੱਤੇ ਸੋਹਣੇ \VUgras{ਨਮੂਨੇ} ਟੰਗੇ ਹਨ ਤੇ
ਛੱਤ ਤੋਂ ਇਕ \VUgras{ਛਾਂ-ਢੱਕਣ} ਸਮੇਤ \VUgras{ਲੈਂਪ} \textsuperscript{(62)} ਲਟਕਿਆ
ਹੈ। \VUgras{ਅਧਿਆਪਕ} \textsuperscript{(60)} ਬਹੁਤ ਧੀਰਜਵਾਨ ਹਨ; ਉਹ ਵਿਦਿਆਰਥੀਆਂ ਦੇ
\VUgras{ਚਿੱਤਰ} \textsuperscript{(61)} ਸੁਧਾਰਦੇ ਹਨ ਤੇ ਉਨ੍ਹਾਂ ਨੂੰ ਸਾਹਮਣੇ ਰੱਖੀ
\VUgras{ਮੂਰਤੀ} \textsuperscript{(59)} ਵੇਖ ਕੇ ਚਿੱਤਰ ਬਣਾਉਣਾ ਸਿਖਾਉਂਦੇ ਹਨ।

%%K t01-16-1 p
16. --- \VUgras{ਸੰਗੀਤ-ਕਮਰਾ} ਬਹੁਤ ਦਿਲਚਸਪ ਲਗਦਾ ਹੈ। ਉੱਥੇ \VUgras{ਸੰਗੀਤ} ਪੜ੍ਹਨਾ,
\VUgras{ਸੁਰ-ਲਿਪੀ} ਉੱਤੇ ਲਿਖੇ \VUgras{ਸੁਰ} \textsuperscript{(67)} ਗਾਉਣਾ (ਡੋ, ਰੇ,
ਮੀ, ਫ਼ਾ, ਸੋਲ, ਲਾ, ਸੀ\,=\,C. D. E. F. G. A. B.), \VUgras{ਕੁੰਜੀਆਂ} ਪਛਾਣਨਾ ਤੇ
ਮਨਮੋਹਣੀਆਂ \VUgras{ਧੁਨਾਂ} ਗਾਉਣਾ ਸਿਖਾਇਆ ਜਾਂਦਾ ਹੈ। ਸੁਰ-ਲਿਪੀਆਂ \VUgras{ਰਹਿਲ}
\textsuperscript{(65)} ਉੱਤੇ ਰੱਖੀਆਂ ਜਾਂਦੀਆਂ ਹਨ। ਇਕ ਵਿਦਿਆਰਥੀ
\VUgras{ਪਿਆਨੋ ਵਜਾ ਰਿਹਾ} \textsuperscript{(64)} ਹੈ ਤੇ \VUgras{ਸੰਗੀਤ-ਅਧਿਆਪਕ}
\textsuperscript{(66)} \VUgras{ਤਾਲ ਦੇ ਰਹੇ} \textsuperscript{(68)} ਹਨ।

%%K t01-17-1 p
17. --- ਸਾਨੂੰ \VUgras{ਦਸਤਕਾਰੀ} \textsuperscript{(69)} ਦੀ \VUgras{ਕਾਰਜਸ਼ਾਲਾ} ਦਾ
ਇਕ ਕੋਨਾ ਮੁਸ਼ਕਲ ਨਾਲ ਵਿਖਾਈ ਦਿੰਦਾ ਹੈ, ਜਿੱਥੇ ਬੱਚੇ ਲੱਕੜ ਛਿੱਲਣਾ ਤੇ ਲੋਹਾ ਰੇਤਣਾ ਸਿੱਖ
ਸਕਦੇ ਹਨ। ਇਹ ਹਰ ਸਕੂਲ ਵਿਚ ਨਹੀਂ ਹੁੰਦੀ।

%%K t01-tit-7 sub
\VUsaut{5.0mm}
\VUpk{10.2pt}{40}{ਵਿਗਿਆਨ-ਕਮਰਾ।}
\VUsaut{3.6mm}

%%K t01-18-1 p
18. --- ਗਣਿਤ ਦੇ ਅਧਿਆਪਕ ਵਿਦਿਆਰਥੀਆਂ ਨੂੰ \VUgras{ਰੇਖਾ-ਗਣਿਤ, ਅੰਕ-ਗਣਿਤ, ਗਿਣਤੀ} ਤੇ
\VUgras{ਮੀਟਰੀ ਪ੍ਰਣਾਲੀ} ਪੜ੍ਹਾਉਂਦੇ ਹਨ। ਸਾਰਣੀ ਉੱਤੇ ਸਾਨੂੰ ਰੇਖਾ-ਗਣਿਤ ਦੀਆਂ ਮੁੱਖ
ਸ਼ਕਲਾਂ ਵਿਖਾਈ ਦਿੰਦੀਆਂ ਹਨ : ਇਕ \VUgras{ਚੱਕਰ} \textsuperscript{(a)}, ਇਕ
\VUgras{ਵਰਗ} \textsuperscript{(b)}, ਇਕ \VUgras{ਤਿਕੋਣ} \textsuperscript{(c)}, ਇਕ
\VUgras{ਰੇਖਾ} \textsuperscript{(d)}।

%%K t01-18-2 p
ਸਾਨੂੰ ਇਕ \VUgras{ਹਿਸਾਬ} ਵੀ ਵਿਖਾਈ ਦਿੰਦਾ ਹੈ; ਉਹ ਨਾ \VUgras{ਜੋੜ} ਹੈ, ਨਾ
\VUgras{ਘਟਾਓ}, ਨਾ \VUgras{ਗੁਣਾ} : ਉਹ \VUgras{ਭਾਗ} \textsuperscript{(e)} ਹੈ।

%%K t01-19-1 p
19. --- ਮੀਟਰੀ ਪ੍ਰਣਾਲੀ ਦੀ ਸਾਰਣੀ ਉੱਤੇ ਸਾਨੂੰ ਵਿਖਾਈ ਦਿੰਦੇ ਹਨ : ਇਕ \VUgras{ਮੀਟਰ}
\textsuperscript{(a)}, ਇਕ \VUgras{ਤੱਕੜੀ} \textsuperscript{(b)} ਤੇ ਉਸ ਦੇ
\VUgras{ਵੱਟੇ}, ਇਕ \VUgras{ਲਿਟਰ} \textsuperscript{(e)}, ਇਕ \VUgras{ਡੈਕਾਲਿਟਰ}
\textsuperscript{(f)}, ਇਕ \VUgras{ਡੈਸੀਲਿਟਰ} ਤੇ ਇਕ \VUgras{ਘਣ ਮੀਟਰ}
\textsuperscript{(d)}।

%%K t01-19-2 p
ਵਿਦਿਆਰਥੀ \VUgras{ਕੁਦਰਤੀ ਵਿਗਿਆਨ} \textsuperscript{(11)} --- ਜੀਵ-ਵਿਗਿਆਨ
\textsuperscript{(a)}, ਬਨਸਪਤੀ-ਵਿਗਿਆਨ \textsuperscript{(b)}, ਭੂ-ਵਿਗਿਆਨ
\textsuperscript{(c)} --- ਤੇ \VUgras{ਇਤਿਹਾਸ} \textsuperscript{(12)} ਵੀ ਪੜ੍ਹਦੇ ਹਨ।

%%K t01-19-3 p
ਅਲਮਾਰੀ ਵਿਚ \VUgras{ਭੌਤਿਕੀ} \textsuperscript{(15)} ਤੇ \VUgras{ਰਸਾਇਣ}
\textsuperscript{(16)} ਦੇ ਸਾਧਨ ਰੱਖੇ ਹਨ।

%%K t01-19-4 p
ਅਲਮਾਰੀ ਦੇ ਉੱਤੇ ਹਨ : ਇਕ \VUgras{ਗੋਲਾ} \textsuperscript{(14)}, ਇਕ \VUgras{ਸ਼ੰਕੂ} ਤੇ
ਇਕ \VUgras{ਭੂ-ਗੋਲਾ} \textsuperscript{(13)}।

%%K t01-19-5 p
ਸਾਰਣੀਆਂ ਦੇ ਉੱਤੇ ਇਕ \VUgras{ਨਕਸ਼ਾ} ਟੰਗਿਆ ਹੈ ਜਿਸ ਉੱਤੇ ਸੰਸਾਰ ਦੇ ਪੰਜੇ ਹਿੱਸੇ ਵਿਖਾਏ
ਗਏ ਹਨ : \VUgras{ਯੂਰਪ} \textsuperscript{(g)}, \VUgras{ਏਸ਼ੀਆ}
\textsuperscript{(e)}, \VUgras{ਅਫ਼ਰੀਕਾ} \textsuperscript{(f)}, \VUgras{ਅਮਰੀਕਾ}
\textsuperscript{(ab)} ਤੇ \VUgras{ਓਸ਼ੀਆਨੀਆ} \textsuperscript{(h)}, ਤੇ ਦੋ ਵੱਡੇ
ਮਹਾਂਸਾਗਰ, \VUgras{ਪ੍ਰਸ਼ਾਂਤ} \textsuperscript{(c)} ਤੇ \VUgras{ਅਟਲਾਂਟਿਕ}
\textsuperscript{(d)}।
\VUsaut{6.0mm}
\VUfilet{22mm}
